% Part: methods
% Chapter: proofs
% Section: reading-proofs

\documentclass[../../../include/open-logic-section]{subfiles}

\begin{document}

\olfileid[id]{mth}{prf}{rea}
\olsection{Membaca Bukti}

Bukti-bukti yang Anda temukan dalam buku teks dan artikel sangat
jarang memberikan semua perincian yang sejauh ini kita sertakan dalam
contoh. Penulis sering tidak menarik perhatian pada saat mereka
membedakan kasus, memberikan bukti tidak langsung, atau menggunakan
suatu definisi. Jadi, ketika membaca bukti dalam buku teks, Anda sering
harus melengkapi sendiri perincian tersebut agar dapat memahami bukti.
Melakukan hal ini juga merupakan latihan yang baik untuk menguasai
berbagai langkah yang harus Anda lakukan dalam suatu bukti. Mari kita
lihat sebuah contoh.

\begin{prop}[Absorpsi]
Untuk semua himpunan $A$, $B$,
\[
A \cap (A \cup B) = A
\]
\end{prop}

\begin{proof}
Jika $z \in A \cap (A \cup B)$, maka $z \in A$, sehingga $A \cap (A
\cup B) \subseteq A$. Sekarang andaikan $z \in A$. Maka $z \in A \cup
B$ juga, dan oleh karena itu $z \in A \cap (A \cup B)$ juga.
\end{proof}

Bukti hukum absorpsi di atas sangat padat. Tidak ada penyebutan
definisi yang digunakan, tidak ada ``kita harus membuktikan bahwa''
sebelum kita membuktikannya, dan seterusnya. Mari kita uraikan bukti
itu. Proposisi yang dibuktikan merupakan klaim umum tentang sebarang
himpunan $A$ dan $B$, dan ketika bukti menyebut $A$ atau $B$, keduanya
merupakan variabel bagi himpunan sembarang. Klaim-klaim umum yang
ditetapkan bukti adalah apa yang diperlukan untuk membuktikan identitas
himpunan, yaitu bahwa setiap !!{element} dari ruas kiri identitas
merupakan !!a{element} dari ruas kanan, dan sebaliknya.

\begin{quote}
``Jika $z \in A \cap (A \cup B)$, maka $z \in A$, sehingga $A \cap (A
  \cup B) \subseteq A$.''
\end{quote}

Ini adalah paruh pertama pembuktian identitas: bagian ini menetapkan
bahwa jika suatu~$z$ sembarang merupakan !!a{element} dari ruas kiri,
maka ia juga merupakan !!a{element} dari ruas kanan, yaitu $A \cap (A
\cup B) \subseteq A$. Asumsikan bahwa $z \in A \cap (A \cup B)$.
Karena $z$ merupakan !!a{element} dari irisan dua himpunan jika dan
hanya jika ia merupakan !!a{element} dari kedua himpunan tersebut,
kita dapat menyimpulkan bahwa $z \in A$ dan juga $z \in A \cup B$.
Khususnya, $z \in A$, dan itulah yang ingin kita tunjukkan. Karena hanya
itulah yang harus dilakukan untuk paruh pertama, kita mengetahui bahwa
sisa bukti pastilah merupakan pembuktian paruh kedua, yaitu bukti bahwa
$A \subseteq A \cap (A \cup B)$.

\begin{quote}
``Sekarang andaikan $z \in A$. Maka $z \in A \cup B$ juga, dan oleh
karena itu $z \in A \cap (A \cup B)$ juga.''
\end{quote}

Kita mulai dengan mengasumsikan bahwa $z \in A$, karena kita sedang
menunjukkan bahwa untuk sebarang~$z$, jika $z \in A$, maka $z \in A
\cap (A \cup B)$. Untuk menunjukkan bahwa $z \in A \cap (A \cup B)$,
kita harus menunjukkan (menurut definisi ``$\cap$'') bahwa (i) $z \in
A$ dan juga (ii) $z \in A \cup B$. Di sini, (i) hanyalah asumsi kita,
sehingga tidak ada lagi yang perlu dibuktikan, dan itulah sebabnya
bukti tidak menyebutnya lagi. Untuk (ii), ingat bahwa $z$ merupakan
!!a{element} dari gabungan himpunan jika dan hanya jika ia merupakan
!!a{element} dari sekurang-kurangnya salah satu himpunan tersebut.
Karena $z \in A$, dan $A \cup B$ adalah gabungan $A$ dan $B$, hal ini
berlaku di sini. Jadi, $z \in A \cup B$. Kita telah menunjukkan baik
(i) $z \in A$ maupun (ii) $z \in A \cup B$; maka, menurut definisi
``$\cap$'', $z \in A \cap (A \cup B)$. Bukti tersebut tidak
menyebutkan definisi-definisi itu; pembaca diasumsikan telah
menginternalisasikannya. Jika Anda belum melakukannya, Anda harus
kembali dan mengingatkan diri tentang definisi-definisi tersebut.
Kemudian Anda juga harus mengenali mengapa dari $z \in A$ mengikuti $z
\in A \cup B$, dan dari $z \in A$ serta $z \in A \cup B$ mengikuti
$z \in A \cap (A \cup B)$.

Berikut versi lain dari bukti di atas, dengan segala sesuatu dinyatakan
secara eksplisit:
\begin{proof}{}
[Menurut definisi $=$ untuk himpunan, untuk menunjukkan $A \cap (A
  \cup B) = A$ kita harus menunjukkan (a) $A \cap (A \cup B)
  \subseteq A$ dan (b) $A \subseteq A \cap (A \cup B)$. (a): Menurut
  definisi $\subseteq$, kita harus menunjukkan bahwa jika $z \in A
  \cap (A \cup B)$, maka $z \in A$.] Jika $z \in A \cap (A \cup B)$,
maka $z \in A$ [sebab menurut definisi $\cap$, $z \in A \cap (A \cup
B)$ jika dan hanya jika $z \in A$ dan $z \in A \cup B$], sehingga $A
\cap (A \cup B) \subseteq A$. [(b): Menurut definisi $\subseteq$, kita
harus menunjukkan bahwa jika $z \in A$, maka $z \in A \cap (A \cup
B)$.] Sekarang andaikan [(1)] $z \in A$. Maka [(2)] $z \in A \cup B$
juga [sebab berdasarkan (1), $z \in A$ atau $z \in B$, yang menurut
definisi $\cup$ berarti $z \in A \cup B$], dan oleh karena itu $z \in
A \cap (A \cup B)$ juga [sebab definisi $\cap$ mensyaratkan $z \in A$,
yaitu (1), dan $z \in A \cup B$, yaitu (2)].
\end{proof}

\begin{prob}
Uraikan bukti $A \cup (A \cap B) = A$ berikut dengan menyebutkan semua
pola inferensi yang digunakan, mengapa setiap langkah mengikuti asumsi
atau klaim yang telah ditetapkan sebelumnya, dan di mana kita harus
merujuk pada definisi tertentu.
\begin{proof}
  Jika $z \in A \cup (A \cap B)$, maka $z \in A$ atau $z \in A \cap
  B$. Jika $z \in A \cap B$, maka $z \in A$. Setiap $z \in A$ juga
  memenuhi $\in A \cup (A \cap B)$.
\end{proof}
\end{prob}

\end{document}
